(1) 理想气体大气的比气体常数为 \(R_s\)、比热为 \(c_p,c_v\)，满足 \(c_p-c_v=R_s\)。竖直坐标 \(z\) 向上、重力加速度 \(g>0\)。若大气处于静水压平衡且沿竖直方向可逆绝热，推导温度直减率。
(2) 给定地面温度 \(T_0\)、压强 \(p_0\)，求温度 \(T(z)\) 与压强 \(p(z)\)，注明物理适用高度范围。
(3) 将一小气团从背景高度平衡位置绝热微位移，推导一般背景 \(T_{\rm env}(z)\) 下的布伦特—魏萨拉频率平方 \(N^2\)。
(4) 分别判断背景温度为等温、干绝热和下降更快时的线性浮力稳定性，并解释环境与气团温度梯度的区别。
